\documentclass[11pt]{article}
\usepackage[margin=25mm]{geometry}
\usepackage{fontspec,unicode-math}
\setmainfont{Latin Modern Roman}
\setmathfont{Latin Modern Math}
\usepackage{amsmath,amssymb,mathtools}
\usepackage{xurl,hyperref}
\hypersetup{colorlinks=true,urlcolor=blue,linkcolor=blue}
\setcounter{secnumdepth}{0}
\setlength{\parindent}{0pt}
\setlength{\parskip}{5pt}
\setlength{\emergencystretch}{3em}
\newcommand{\sref}[2]{\hyperlink{source:#1:#2}{[S#2]}}
\newcommand{\cataloguescope}[1]{\paragraph{Source question.}#1}
\newcommand{\floor}[1]{\left\lfloor#1\right\rfloor}
\begin{document}
\section{Decidability of Simplicial-Complex Embeddability}\label{statement}

\subsection{Definitions and mathematical statement}

\paragraph{Shared notation.}Standard number systems are denoted by $\mathbb N,\mathbb Z,\mathbb Q,\mathbb R,\mathbb C$; index ranges, computational models and specialized notation follow the cited problem source and the definitions below.


A polyhedron is an intersection of finitely many closed affine half-spaces; a polytope is a bounded polyhedron. Its faces are its intersections with supporting hyperplanes.

A finite simplicial complex is a finite family of vertex subsets closed under taking subsets. Its dimension is one less than its largest simplex size. A piecewise-linear embedding is an injective map linear on simplices after finite subdivisions.

A continuous piecewise-linear function is continuous and agrees with an affine function on each member of a finite polyhedral partition. A difference-of-convex decomposition writes it as $g-h$ with both $g$ and $h$ convex.
\paragraph{Statement recorded in the supplied source.}
For fixed integers $k\leq d$, is there an algorithm that decides whether the polyhedron of a given finite simplicial complex $K$ of dimension at most $k$ admits a piecewise-linear embedding into $\mathbb R^d$? \sref{statement}{2}
\paragraph{Scope qualification.}
The target is a classification over fixed dimension pairs, not a blanket claim that one algorithm must work uniformly in all dimensions.
\paragraph{Source status.}
The supplied catalogue labels this record partially solved. Its source-status metadata was checked against the supplied snapshot on 2026-10-11. This import does not independently establish current open status or mathematical correctness. The catalogue literature assessment is dated 2026-08-21.
\subsection{Short English statement}
Determine decidability of piecewise-linear embedding for all the stated source and ambient dimensions.
\subsection{Sources}
\begin{itemize}
\item[\hypertarget{source:statement:1}{S1}] Ulam AI and the UnsolvedMath Contributors, UnsolvedMath: A Curated Collection of Open Mathematics Problems, catalogue entry. Catalogue curation and metadata: CC BY 4.0. The supplied source question is retained; definitions and the English summary are editorial additions.
\url{https://huggingface.co/datasets/ulamai/UnsolvedMath}.
\item[\hypertarget{source:statement:2}{S2}] Discrete Geometry (2009), Oberwolfach Reports, DOI 10.4171/owr/2008/44, pp. 2477--2556
\url{https://doi.org/10.4171/owr/2008/44}.
\item[\hypertarget{source:statement:3}{S3}] Matousek--Tancer--Wagner
\url{https://arxiv.org/abs/0807.0336}.
\textit{Underlying Oberwolfach Reports material retains its source-specific licence, including CC BY-SA 4.0 where applicable; see THIRD\_PARTY\_NOTICES.md.}
\end{itemize}
\paragraph{Preparation and review.}
Original question attributed to the cited source and UnsolvedMath contributors. Codex (GPT-6; exact serving revision not exposed) assisted with definitions, English summaries, source inspection and formatting. No new human mathematical authorship or independent proof verification is asserted. Catalogue identity checked against the complete supplied registry and existing repository statements; known flagged repetitions and obvious identical conjectures excluded. Full mathematical-equivalence/literature review has not been performed. Statement preparation and final review followed contribution prompts C and D; automated checks verify structure and identity, not open status.
\label{end:statement}
\end{document}
